% Candidate transfer of existing local results, not a claim of a new theorem.
% Sources: trianglefree_list_recolouring_{path_factor,cut_bound}_20260915_zh.md.
% The Akiyama--Avis--Era statement was checked in Egawa--Furuya--Ozeki,
% https://tgt.ynu.ac.jp/ozeki/2015EFO.pdf, Theorem A, PDF page 2.
% Ellingham--Zha Theorem 3.5 was checked in the authors' manuscript:
% https://math.vanderbilt.edu/ellingmn/paper/ttw/b-eps.pdf, PDF pages 9--11.

\begin{lemma}[Analytic path factors and bounded-load connections]
Let $X$ be the complete fixed-list recolouring graph of a finite
triangle-free graph. Every nontrivial component $Y$ of $X$ has a
$\{P_2,P_3\}$-factor. Given any such factor $F$, its components can
be joined by $c(F)-1$ edges of $Y$, with each physical vertex of $Y$
incident with at most three of the added edges.
\end{lemma}
\begin{proof}
We first show $c(Y-S)\le2|S|$ for every nonempty $S\subseteq V(Y)$.
For a component $D$ of $Y-S$ put
$b(D)=|N_Y(D)\cap S|$, which is positive by connectedness.
At each $s\in S$, choose one neighbour $i_D\in D$ for every
component adjacent to $s$. These representative moves change
different original vertices: moves at the same vertex give adjacent
states and hence cannot lead to different components of $Y-S$.

Join two representatives in an auxiliary graph $C_s$ when their
moves can be performed simultaneously. Two distinct-coordinate
legal moves fail to commute precisely when their original vertices
are adjacent and their target colours agree. Thus $\alpha(C_s)\le2$,
since three pairwise noncommuting moves would create an original
triangle. For a fixed $D$, each commuting partner produces a
distinct double-move state adjacent to $i_D$. It must belong to $S$;
otherwise it joins two components of $Y-S$. It is different from $s$,
so $d_{C_s}(D)+1\le b(D)$.

For any finite graph $C$, a uniform random ordering and selection
of the vertices earliest in their closed neighbourhoods give
$\sum_{v\in V(C)}1/(d_C(v)+1)\le\alpha(C)$. Consequently
\[
\sum_{D:s\in N_Y(D)}\frac1{b(D)}
\le\sum_{D\in V(C_s)}\frac1{d_{C_s}(D)+1}\le2.
\]
Summing over $s\in S$ counts each component $D$ once, proving the
cut bound. It implies $i(Y-S)\le2|S|$ for every $S$, including
$S=\varnothing$ because $Y$ is nontrivial and connected. The
Akiyama--Avis--Era path-factor criterion therefore gives a
$\{P_2,P_3\}$-factor. Finally, the cut bound gives
$c(Y-S)<2|S|+3$ also for $S=\varnothing$. Ellingham--Zha,
Theorem~3.5 with $g=2,k=3$, applied to the specified factor $F$,
gives exactly the asserted connections and physical load.
\end{proof}

\begin{corollary}[Listed pentagons]
For lists of size at least three in a four-colour palette, every
component of the complete listed $C_5$ recolouring graph satisfies
the preceding lemma.
\end{corollary}
\begin{proof}
Every proper colouring repeats a colour at two of the five vertices.
The repeated vertices are nonadjacent, hence have a common neighbour
on $C_5$. At that common neighbour only one colour is forbidden by
its two neighbours; its list still contains at least two choices.
Thus every state has a legal move and no component is a singleton.
\end{proof}

% Proposed replacement for the protected-product proof, preserving K_A,
% every existing cycle/seed, every pool, and all 313 interface witnesses.
\begin{proof}[Candidate proof of protected-product]
Choose a $\{P_2,P_3\}$-factor of $B$. For each factor path use a
cycle of $A\square P_2$ or $A\square P_3$: if $A$ is Hamiltonian,
splice its two or three copies using the first two pool edges;
if $A$ is hard, use the existing two- or three-row seed directly.
These cycles retain $K_A$ and the last six pool edges in every row.

Connect the factor paths using the preceding lemma, with at most
three connecting edges incident at any physical $B$ vertex. Fix any
three of the last six pool edges as splice edges. Process the tree
of factor paths in parent-before-child order. At a parent row at most
two of the three splice edges have been consumed; the fresh child
row retains all three. Choose an unused common splice edge and
perform the ordinary two-edge switch. It joins two disjoint cycles
into one and leaves every edge of $K_A$ intact. Induction covers all
factor paths, giving the required protected Hamilton cycle.
\end{proof}

% Consequences and limits:
% 1. The 590-word factor table is no longer a premise of this proof.
% 2. Arbitrary-r strip construction and hard seed-block concatenation vanish.
% 3. All 23 component cycles/seeds, K_A definitions and 313 interfaces remain.
% 4. Factor 00000's two C15 words are still used as component cycles.
% 5. Ordinary (non-factor) cycles and the 00112 seeds must be retained.
% 6. This is an analytic transfer of earlier local material and classical
%    external results; it is neither formal verification nor a new claim
%    resolving the unrestricted planar girth-five Hamiltonicity problem.
